\documentclass{article}
\usepackage[utf8]{inputenc}
\usepackage{graphicx}
\usepackage{amsmath}
\usepackage{booktabs}
\usepackage{hyperref}
\usepackage{float}

\title{Report 2: Fetal Head Circumference Measurement using Ultrasound}
\author{Nguyen Bich Thu}
\date{}

\begin{document}
\maketitle

\begin{abstract}
This report presents a deep learning approach for measuring fetal head circumference (HC) from 2D ultrasound images using the HC18 dataset. A dual-input regression model combining MobileNetV3 features with pixel size information is implemented. The model achieves a training MAE of 16.41 mm after 10 epochs, demonstrating the effectiveness of incorporating physical pixel size information into the regression task.
\end{abstract}

\section{Introduction}
Fetal head circumference (HC) measurement is a crucial biometric parameter in prenatal care for monitoring fetal growth and detecting potential abnormalities. Traditionally, this measurement requires manual annotation by sonographers, which is time-consuming and subject to inter-observer variability. This project aims to automate HC measurement using deep learning techniques on ultrasound images from the HC18 dataset.

\section{Dataset Description}
\subsection{Dataset Overview}
The HC18 dataset \cite{van2018hc18} is a publicly available dataset for fetal head circumference measurement. It consists of:
\begin{itemize}
    \item \textbf{Training set}: 999 2D ultrasound images with corresponding pixel size and HC measurements
    \item \textbf{Testing set}: 335 2D ultrasound images (without public ground truth labels)
\end{itemize}

\subsection{Data Characteristics}
\begin{itemize}
    \item \textbf{Image format}: PNG files, grayscale ultrasound images
    \item \textbf{Image dimensions}: Varying sizes (typically around 540$\times$800 pixels)
    \item \textbf{Pixel size range}: 0.052 mm to 0.326 mm per pixel
    \item \textbf{HC range}: 45.19 mm to 371.96 mm (covering various gestational ages)
\end{itemize}

\subsection{Data Distribution}
The dataset includes CSV files containing metadata:
\begin{itemize}
    \item Filename: Image identifier
    \item Pixel size (mm): Physical size of each pixel
    \item Head circumference (mm): Ground truth measurement
\end{itemize}

\section{Methodology}
\subsection{Model Architecture}
A dual-input regression model was implemented combining image features with pixel size information:

\begin{itemize}
    \item \textbf{Backbone}: MobileNetV3-Small (pretrained on ImageNet) for feature extraction
    \item \textbf{Feature dimension}: 576 features from the backbone
    \item \textbf{Regression head}: 
    \begin{itemize}
        \item Input: 576 image features + 1 pixel size value
        \item Hidden layers: [577 $\rightarrow$ 128 $\rightarrow$ 64 $\rightarrow$ 1]
        \item Activation: ReLU
        \item Dropout: 0.2 for regularization
    \end{itemize}
\end{itemize}

The model architecture can be represented as:
\begin{equation}
    \hat{HC} = f(g(I) \oplus p)
\end{equation}
where $g(I)$ represents image features extracted by MobileNetV3, $p$ is the pixel size, and $\oplus$ denotes concatenation.

\subsection{Data Preprocessing}
Images were preprocessed using the following steps:
\begin{enumerate}
    \item Resize to 224$\times$224 pixels (standard input size for ImageNet models)
    \item Convert to RGB format (3 channels)
    \item Normalize using ImageNet statistics:
    \begin{itemize}
        \item Mean: [0.485, 0.456, 0.406]
        \item Std: [0.229, 0.224, 0.225]
    \end{itemize}
\end{enumerate}

\subsection{Training Configuration}
\begin{itemize}
    \item \textbf{Loss function}: Mean Absolute Error (L1 loss)
    \item \textbf{Optimizer}: Adam (learning rate = 0.001)
    \item \textbf{Batch size}: 8
    \item \textbf{Epochs}: 10
    \item \textbf{Device}: CPU (can be adapted for GPU)
\end{itemize}

\subsection{Implementation Details}
The model was implemented using PyTorch with the following key components:
\begin{itemize}
    \item Custom Dataset class for loading images and metadata
    \item Dual-input architecture combining CNN features with scalar input
    \item Transfer learning using pretrained MobileNetV3 weights
\end{itemize}

\section{Results and Discussion}
\subsection{Training Progress}
The model was trained for 10 epochs on the training set (999 images). Figure \ref{fig:training} shows the training MAE over epochs, demonstrating consistent improvement from 61.83 mm to 16.41 mm.

\begin{figure}[H]
    \centering
    \includegraphics[width=0.8\textwidth]{image.png}
    \caption{Training Progress: Mean Absolute Error over 10 epochs}
    \label{fig:training}
\end{figure}

\begin{table}[H]
    \centering
    \caption{Training MAE Progression}
    \begin{tabular}{cc}
        \toprule
        Epoch & Training MAE (mm) \\
        \midrule
        1 & 61.83 \\
        2 & 27.49 \\
        3 & 25.00 \\
        4 & 23.05 \\
        5 & 23.18 \\
        6 & 20.40 \\
        7 & 19.22 \\
        8 & 18.25 \\
        9 & 17.59 \\
        10 & 16.41 \\
        \bottomrule
    \end{tabular}
\end{table}

\subsection{Comparison with Leaderboard}
The current implementation achieves a training MAE of 16.41 mm after 10 epochs. For context, the HC18 challenge leaderboard results \cite{hc18leaderboard} show:

\begin{table}[H]
    \centering
    \caption{Comparison with HC18 Leaderboard (Test Set Performance)}
    \begin{tabular}{lcc}
        \toprule
        Method & MAE (mm) & DSC (\%) \\
        \midrule
        Top-performing methods & 2.0-3.0 & 97-98 \\
        Our model (training) & 16.41 & - \\
        \bottomrule
    \end{tabular}
\end{table}

\subsection{Hyperparameter Experiments}
Several hyperparameter configurations were explored:

\begin{table}[H]
    \centering
    \caption{Hyperparameter Experiments}
    \begin{tabular}{lcc}
        \toprule
        Configuration & Final MAE (mm) & Notes \\
        \midrule
        Baseline (lr=0.001, batch=8) & 16.41 & Current model \\
        Higher learning rate (0.01) & \textgreater 100 & Unstable training \\
        Smaller batch size (4) & 19.23 & Slower convergence \\
        No dropout & 18.75 & Slight overfitting \\
        \bottomrule
    \end{tabular}
\end{table}

\section{Conclusion}
This project demonstrates a deep learning approach for automated fetal head circumference measurement. Key achievements include:
\begin{itemize}
    \item Successful implementation of a dual-input regression model
    \item Integration of image features with physical pixel size information
    \item Achievement of 16.41 mm MAE on training data
\end{itemize}

\begin{thebibliography}{9}
\bibitem{van2018hc18}
van den Heuvel, T. L., de Bruijn, D., de Korte, C. L., \& Ginneken, B. V. (2018).
\newblock Automated measurement of fetal head circumference using 2D ultrasound images.
\newblock \textit{PLOS ONE}, 13(8), e0200412.

\bibitem{hc18leaderboard}
HC18 Challenge Leaderboard.
\newblock \url{https://hc18.grand-challenge.org/evaluation/challenge/leaderboard/}
\end{thebibliography}

\end{document}